% Part: first-order-logic
% Chapter: proof-systems
% Section: natural-deduction

\documentclass[../../../include/open-logic-section]{subfiles}

\begin{document}

\iftag{FOL}
      {\olfileid{fol}{prf}{ntd}}
      {\olfileid{pl}{prf}{ntd}}

\olsection{Natuurlijke deductie}

Natuurlijke deductie is !!a{derivation}ssysteem dat werkelijke
redeneringen beoogt te weerspiegelen, in het bijzonder de
gereglementeerde redeneringen van wiskundigen. Werkelijke redeneringen
volgen een aantal ``natuurlijke'' patronen. Met een bewijs door
gevalsonderscheid kan uit een disjunctieve premisse een conclusie worden
afgeleid door aan te tonen dat de conclusie uit elk van de disjuncten
volgt. Met een bewijs uit het ongerijmde kan een conclusie worden
afgeleid door aan te tonen dat haar negatie tot een tegenspraak leidt.
Een implicatie ``als \dots, dan \dots'' kan worden bewezen door aan te
tonen dat het consequens uit het antecedens volgt. Natuurlijke deductie
formaliseert enkele van deze natuurlijke afleidingspatronen. Bij elk
logisch connectief en elke kwantor horen twee regels: een
introductieregel en een eliminatieregel. Elke regel komt overeen met
zo'n natuurlijk afleidingspatroon. Zo komt $\Intro{\lif}$ overeen met
het bewijzen van een implicatie en $\Elim{\lor}$ met een bewijs door
gevalsonderscheid. Een bijzonder eenvoudige regel is $\Elim{\land}$:
daarmee kan uit $!A \land !B$ de formule~$!A$ (of $!B$) worden afgeleid.

Natuurlijke deductie onderscheidt zich van andere
!!{derivation}ssystemen onder meer door het gebruik van aannames.
!!^a{derivation} in natuurlijke deductie is een boom van
!!{formula}s. Aan de wortel van de boom staat één !!{formula}; de
``bladeren'' zijn de !!{formula}s waaruit de conclusie wordt afgeleid.
Sommige !!{formula}s aan de bladeren spelen binnen de
!!{derivation} een rol, maar zijn ``opgebruikt'' wanneer de conclusie
wordt bereikt. Dit komt overeen met de gangbare praktijk om bij een
redenering hypothesen in te voeren die slechts korte tijd van kracht
blijven. Bij een bewijs door gevalsonderscheid nemen we bijvoorbeeld
elk van de disjuncten aan; bij het bewijzen van een implicatie nemen we
het antecedens aan; bij een bewijs uit het ongerijmde nemen we de
negatie van de conclusie aan. Kenmerkend voor natuurlijke deductie is
dat zulke hypothetische aannames worden ingevoerd en vervolgens, om
een tussenstap vast te stellen, weer worden opgeheven. De !!{formula}s
aan de bladeren van !!a{derivation} heten aannames; sommige
afleidingsregels kunnen ze ``!!{discharge}.'' Als we bijvoorbeeld
!!a{derivation} van~$!B$ hebben uit aannames waaronder~$!A$, dan mogen
we volgens $\Intro{\lif}$ de formule~$!A \lif !B$ afleiden en elke
aanname van de vorm~$!A$ !!{discharge}. (Om bij te houden welke
aannames bij welke afleidingsstappen worden opgeheven, geven we de
afleidingsstap en de aannames die daarbij worden opgeheven hetzelfde
nummer.) De aannames die aan het eind van de !!{derivation}
!!{undischarged} zijn, zijn samen voldoende voor de waarheid van de
conclusie. !!^a{derivation} stelt dus vast dat de conclusie volgt uit
de !!{undischarged} aannames.

De op natuurlijke deductie gebaseerde relatie $\Gamma \Proves !A$
geldt dan en slechts dan als er !!a{derivation} bestaat waarin $!A$~de
laatste !!{sentence} van de boom is en elk !!{undischarged} blad
tot~$\Gamma$ behoort. $!A$~is een stelling in natuurlijke deductie dan
en slechts dan als er !!a{derivation} bestaat waarin $!A$~de laatste
!!{sentence} is en alle aannames zijn !!{discharged}. De volgende
!!a{derivation} laat bijvoorbeeld zien dat $\Proves (!A \land !B) \lif
!A$:
\begin{prooftree}
  \AxiomC{$\Discharge{!A \land !B}{1}$}
  \RightLabel{\Elim{\land}}
  \UnaryInfC{$!A$}
  \DischargeRule{\Intro{\lif}}{1}
  \UnaryInfC{$(!A \land !B) \lif !A$}
\end{prooftree}
Het nummer~$1$ geeft aan dat de aanname $!A \land !B$ bij de
\Intro{\lif}-afleidingsstap wordt !!{discharged}.

Een verzameling~$\Gamma$ is inconsistent dan en slechts dan als
$\Gamma \Proves \lfalse$ in natuurlijke deductie. De regel \FalseInt{}
zorgt ervoor dat uit een inconsistente verzameling elke !!{sentence}
kan worden !!{derive}d.

Systemen voor natuurlijke deductie zijn in de jaren dertig van de
twintigste eeuw ontwikkeld door Gerhard Gentzen en Stanis\l{}aw
Ja\'skowski en later verder ontwikkeld door Dag Prawitz en Frederic
Fitch. Omdat de afleidingsstappen natuurlijke bewijsmethoden
weerspiegelen, geniet natuurlijke deductie de voorkeur van filosofen.
De door Fitch ontwikkelde versies worden vaak gebruikt in inleidende
leerboeken over logica. Binnen de filosofie van de logica worden de
regels van de natuurlijke deductie soms opgevat als bepalend voor de
betekenis van de logische operatoren (``bewijstheoretische semantiek'').

\end{document}
